\documentclass[../../../include/open-logic-section]{subfiles}

\begin{document}

\olfileid[hi]{sth}{story}{approach}
\olsection{संचयी-पुनरावर्ती दृष्टिकोण}

हम समुच्चयों को चरणों में किस प्रकार स्तरित करेंगे, इसका कुछ अधिक पूर्ण
विवरण यह है:
\begin{quote}
	समुच्चय \emph{चरणों} में बनते हैं। प्रत्येक चरण $S$ के लिए कुछ चरण
	$S$ से \emph{पहले} होते हैं। चरण $S$ पर, $S$ से पहले के चरणों में बने
	समुच्चयों से बने प्रत्येक संग्रह को एक समुच्चय बनाया जाता है। चरणों में
	बने समुच्चयों के अतिरिक्त कोई समुच्चय नहीं हैं।
	\citep[p.~323]{Shoenfield:AST}
\end{quote} 
यह \emph{समुच्चय की संचयी-पुनरावर्ती अवधारणा} की रूपरेखा है। यही
\olref[sth][][]{part} में प्रस्तुत औपचारिक समुच्चय सिद्धांत का आधार होगी।

इसे कुछ और विस्तार से देखें। शोएनफ़ील्ड के वर्णन में प्रत्येक चरण पर हम उन
समुच्चयों से नए समुच्चय बनाते हैं जो पहले के चरणों से उपलब्ध हैं। अतः
शोएनफ़ील्ड के चित्र में आरंभिक चरण, चरण $0$, से \emph{पहले} कोई चरण नहीं
है और इसलिए \emph{और भी अधिक स्पष्टतः} पहले के चरणों से कोई समुच्चय उपलब्ध
नहीं है।\footnote{हम क्यों मानें कि कोई \emph{पहला} चरण है? 
\olref[z][story]{sec} में \stagesord{} की पादटिप्पणी देखें।} अतः हम केवल
एक समुच्चय बनाते हैं: बिना किसी !!{element} वाला समुच्चय $\emptyset$। चरण
$1$ पर पहले के चरणों से ठीक एक समुच्चय उपलब्ध है, इसलिए केवल एक नया समुच्चय
$\{\emptyset\}$ है। चरण $2$ पर पहले के चरणों से दो समुच्चय उपलब्ध हैं और
हम दो नए समुच्चय $\{\{\emptyset\}\}$ तथा
$\{\emptyset, \{\emptyset\}\}$ बनाते हैं। चरण $3$ पर पहले के चरणों से चार
समुच्चय उपलब्ध हैं, इसलिए हम बारह नए समुच्चय बनाते हैं\ldots। इस प्रकार
समुच्चयों का संचयी-पुनरावर्ती चित्र कुछ ऐसा दिखेगा (संख्याएँ चरण दर्शाती
हैं):
\begin{center}
	\begin{tikzpicture}[scale=0.6]
	\tikzset{cut_here/.style={densely dotted}}
	\draw (-4,6) -- (-1, 0)-- (2,6); 
	\draw[cut_here] (-5,8)--(-4,6);
	\draw[cut_here] (2,6)--(3,8);
	\draw(-1.5, 1)--(-0.5, 1);
	\draw(-2, 2)--(0, 2);
	\draw(-2.5, 3)--(0.5,3);
	\draw(-3, 4)--(1, 4);
	\draw(-3.5, 5)--(1.5, 5);
	\draw(-4, 6)--(2, 6);
	\node[label] at (0, 0) {\small 0};
	\node[label] at (0.5, 1) {\small 1};
	\node[label] at (1, 2) {\small 2};
	\node[label] at (1.5, 3) {\small 3};
	\node[label] at (2, 4) {\small 4};
	\node[label] at (2.5, 5) {\small 5};
	\node[label] at (3, 6) {\small 6};
	\end{tikzpicture}
\end{center}
तो इस कथा को क्यों स्वीकार करें?

एक कारण यह है कि यह अच्छी और सुगम कथा है। सबसे स्पष्ट कथा, अर्थात सहज
समुच्चय-निर्माण, के विफल हो जाने के बाद हमें किसी अच्छी कथा की आवश्यकता है।

किंतु यह कथा \emph{केवल} अच्छी नहीं है। यह मानने का अच्छा कारण है कि इस
कथा पर आधारित कोई भी समुच्चय सिद्धांत \emph{संगत} होगा। कारण यह है।

समुच्चय की संचयी-पुनरावर्ती अवधारणा के अनुसार हम चरणों में समुच्चय बनाते हैं
और उनके !!{element}s ऐसी वस्तुएँ होनी चाहिए जो \emph{पहले से} उपलब्ध थीं।
अतः किसी भी चरण~$S$ के लिए हम यह समुच्चय बना सकते हैं:
\[
	R_S = \Setabs{x}{x \notin x \text{ और $x$, $S$ से पहले उपलब्ध था}}
\]
रसेल का विरोधाभास सिद्ध करने वाला तर्क अब स्थापित करेगा कि $R_S$ स्वयं
चरण $S$ से पहले उपलब्ध नहीं है। और यह कोई विरोधाभास नहीं। इसके अतिरिक्त,
यदि हम समुच्चय की संचयी-पुनरावर्ती अवधारणा स्वीकार करते हैं, तो हमें स्वयं
रसेल समुच्चय बना पाने की \emph{अपेक्षा} भी नहीं करनी चाहिए थी। क्योंकि वह
उन सभी स्वयं के असदस्य समुच्चयों का समुच्चय होता जो ``कभी भी उपलब्ध होंगे''।
संक्षेप में: संचयी-पुनरावर्ती कथा को देखते हुए यह तथ्य \emph{आश्चर्यजनक}
नहीं कि हम (सिद्धतः) रसेल समुच्चय नहीं बना सकते; यही तो हमारी
\emph{भविष्यवाणी} होगी।

%In one sense, then, the cumulative-iterative conception of set yields a response to Russell's Paradox which is rather like the \emph{predicativist}'s response. After all, the predicativist said that the collection of all non-self-membered sets$_n$ is a set$_{n+1}$, and this set$_n$ will not be self-membered. (Indeed, most predicativists will treat it as \emph{ungrammatical} to try to ask whether a set is self-membered.)
%
%But there is an important difference between the cumulative-iterative approach and the predicativist approach: the cumulative-iterative approach treats all of our entities as being \emph{of the same kind}. The predicativist will presumably have an empty set$_0$ (i.e., a set$_0$ with no !!{element}s), and an empty set$_1$ (i.e., a set$_1$ with no !!{element}s), and these will be \emph{different entities}. But on the cumulative-iterative approach, there is just one empty set, $\emptyset$, and it will be available at every stage of the hierarchy. 

%Indeed, the Axiom of Extensionality, stated at the start of this chapter, will hold true of the elements in this hierarchy (so that there can be \emph{only one} ``empty'' set). But I do not intend to use the cumulative-iterative conception to justify Extensionality (see \cite{Boolos1971}). Again: I suggest that we should accept Extensionality, just because we are interested in extensional collections. 

\end{document}
